\begin{lemma}[The zero-divider PE predecessor color]
Let $P$ be a Step 1 exit-leg presentation whose good hierarchy has zero
dividers and whose terminal certified extension kind is PE.  Suppose
that the terminal ladder tree is elementary under the coloring
immediately preceding that PE step.  Then there is a color $\theta$,
different from the two exit-leg colors, which is missing and
singleton-closed on the last starred stage and singleton-closed on the
terminal tree of $P$.
\end{lemma}
\begin{proof}
Let $X$ be the terminal ladder tree, let $B$ be the hierarchy base, and
write $A=B.\mathit{core}$ for the zero-divider starred stage.  The
`parallelBase' field of \noderef{GoodHierarchy} gives
\[
 4\leq\left|\bigl(\overline\psi(X)\cap
 \overline\psi(B.\mathit{closure})\bigr)
 \setminus\bigcup_{\eta\in D^K_{n,0}}\Gamma^0_\eta\right|.
\]
After deleting the two distinct exit-leg colors, choose a remaining
color $\theta$.  It is missing on both $X$ and the base closure and
avoids the internal Gamma union.

The assumed pre-step elementariness is precisely the remaining input of
\noderef{ParallelCommonMissingClosed}; the ETT presentation identifies
its generating coloring with the last series coloring.  That theorem
therefore closes every common missing color, including $\theta$, on the
base core.  The `parallel' field of \noderef{HierarchyBase} and the
vertex-set clause of \noderef{OrderedTreeUnion} identify this core with
the exact ordered union of $X$ and the base closure, while
\noderef{HierarchyStarredStage} identifies it with $A$.  Consequently
$\theta$ is missing and singleton-closed on $A$.

The hierarchy inclusions place $X$ inside the terminal tree, so
$\theta$ is missing there.  From the definitions
\noderef{HierarchyTerminalStoredColors} and
\noderef{HierarchyStoredColorsAt}, inclusion of missing-color sets gives
that the terminal stored set is contained in the index-zero stored set.
Thus avoidance of the internal Gamma union implies avoidance of the
terminal Gamma union.  The terminal minus-closure field of
\noderef{Step1ExitLegPresentation}, unfolded through
\noderef{TreeColorSetMinusClosed}, closes $\theta$ on the terminal tree.
\end{proof}
